\section{Métodos conservadores}

\subsection{CQL: Conservative Q-Learning}

La idea de CQL es más directa: agregar un \textbf{término de penalización} en la actualización de la función Q, para reducir activamente los valores Q.

$$\hat{Q} = \arg\min_{Q} \max_\mu \; \mathbb{E}_{(s,a,s') \sim \mathcal{D}}\left[(Q(s,a) - (r + \gamma \mathbb{E}_\pi\left[Q(s', a')\right]))^2\right] + \alpha \mathbb{E}_{s \sim \mathcal{D}, a \sim \mu(\cdot|s)}\left[Q(s,a)\right] - \alpha \mathbb{E}_{(s,a) \sim \mathcal{D}}\left[Q(s,a)\right]$$

\begin{itemize}
    \item Primer término: actualización TD estándar.
    \item Segundo término: reduce el valor Q de todas las acciones ($\mu$ cubre todo el espacio de acciones).
    \item Tercer término: restablece el valor Q de las acciones presentes en los datos (para no reducirlo en exceso).
\end{itemize}

Puede demostrarse que, para un $\alpha$ suficientemente grande, el valor Q aprendido es \textbf{conservador} (subestimado en vez de sobreestimado) bajo la distribución de la política.

\subsection{Implementación de CQL}

Cuando se usa regularización de máxima entropía, el $\mu(a|s) \propto \exp(Q(s,a))$ óptimo hace que el término de penalización se convierta en $\log \sum_a \exp(Q(s,a))$, sin necesidad de construir $\mu$ de manera explícita.

\subsection{Resumen de la sección}

CQL reduce activamente los valores Q para evitar la sobreestimación, y es otro método eficaz de RL fuera de línea. A diferencia de la restricción implícita de IQL, CQL construye una estimación conservadora de manera explícita.

